% Author: Charis Opol (F-11). Draft of 23 Nov 2026, built from the proposal's introduction and
% the repository. Numbers: analysis/01_grid/tables/test_wer_greedy.md, Table~\ref{tab:realworld}
% (reports/flix_v5/final_decoder.md) and Section~\ref{sec:results}.
\section{Introduction}
\label{sec:introduction}

\PARstart{A}{utomatic} speech recognition (ASR) performs well on clean audio, but its error rate
rises quickly with background noise. For live captioning, used by deaf and hard-of-hearing people
in lectures and meetings, these errors make captions least reliable where they are most needed.
Visual speech information is unaffected by acoustic noise. Lip movements alone are ambiguous, since
phonemes such as /p/ and /b/ share a mouth shape, but they complement the audio in exactly the
conditions where the audio fails. Audio-visual speech recognition (AVSR) exploits this, and
published systems report their largest gains at low signal-to-noise ratios (SNRs)
(Section~\ref{sec:related}).

This paper addresses two research questions:
\begin{enumerate}
\item Does lip-assisted fusion reduce recognition errors across noise levels on a controlled
benchmark, the GRID corpus~\cite{cooke2006} with MUSAN noise~\cite{snyder2015} from clean to
$-5$~dB, tested on unseen speakers?
\item Does the gain carry over to open-vocabulary speech, measured on held-out film and TED clips
against an audio-only model that differs only in the absence of the lips?
\end{enumerate}
The project proposal also listed emotion detection and spoken-language identification. These were
removed from the core system in the implementation plan, because GRID provides no emotion labels
and a single language.

The contributions are as follows.
\begin{itemize}
\item Model P4, a model trained with connectionist temporal classification (CTC) that fuses an audio encoder and a ResNet-18 lip encoder through
bidirectional cross-modal attention, a per-frame gate, a learned missing-stream embedding and
modality dropout (Section~\ref{sec:model}). On unseen GRID speakers it roughly halves the word error
rate (WER) of the audio-only baseline at every noise level (Section~\ref{sec:results}).
\item Ablations of modality dropout, attention with gating, and the 3-D front end, each with a
paired bootstrap interval.
\item A real-world model, ws\_v5, that fuses pretrained audio and lip encoders through a
reliability-gated fusion initialised as its audio-only twin (Section~\ref{sec:pipelines-real}).
On held-out TED talks at $-10$~dB it reduces WER from 42.3~\% to 28.5~\% (Section~\ref{sec:realworld}).
\item A negative result: when the audio does not belong to the visible speaker, ws\_v5 transcribes
the wrong voice, and retraining on mismatched audio trades this error for a large loss on clean
speech.
\item An application in which a command-line tool, a Gradio interface, a browser demo and a live
capture mode share one inference engine with the evaluation code (Section~\ref{sec:application}).
\end{itemize}

Section~\ref{sec:related} reviews related work. Sections~\ref{sec:data} to~\ref{sec:model} describe
the data, pipelines and models, and Sections~\ref{sec:application} and~\ref{sec:testing} the
application and its tests. Sections~\ref{sec:results} and~\ref{sec:realworld} report results on
GRID and on real speech. Section~\ref{sec:limitations} discusses limitations, and
Section~\ref{sec:conclusion} concludes.
